% ─────────────────────────────────────────────────────────────────
%  Draft of §IV.D (hidden hot-runaway + IC hysteresis prediction) —
%  to replace existing §IV.E (Hidden multistability, lines 609–696
%  of main.tex).  Pulls in figure D2 (hysteresis_C_theta0) and
%  D1 (refined basin at C-cell), and inherits Fig.6 (a),(b),(c).
% ─────────────────────────────────────────────────────────────────

\subsection{A hidden hot-runaway state and the IC-induced hysteresis
prediction}
\label{sec:multistability}

Sec.~\ref{sec:lsa_vs_nl} closed by asking whether the C-region of
the parameter plane—where the swollen homogeneous SS does not exist
yet the PDE oscillates—corresponds to genuine attractor multiplicity
of the spatially extended problem.  We resolve this by initial-
condition (IC) sweeps at a representative C-cell
($\Bi_T=0.059$, $S_\chi=1.80$).

\paragraph{Two co-existing attractors.}
At identical bath conditions and reaction parameters, simulations
initialized in different regions of $(J_0, \theta_0)$ space converge
to one of two qualitatively distinct attractors
[Fig.~\ref{fig:fig_manifold}(c)]: (i)~the LCST-front cycle of
Sec.~\ref{sec:lcst_front_mechanism}, with bulk-mean
$\langle J\rangle\!\approx\!2$,
$\langle\theta\rangle\!\approx\!2$, and (ii)~a previously-unreported
\emph{hot-runaway} steady state with
$\langle J\rangle\!\approx\!5$,
$\langle\theta\rangle\!\approx\!12$.
The two attractors differ by
$\Delta\langle J\rangle\!\approx\!3.3$ and
$\Delta\langle\theta\rangle\!\approx\!10$, far beyond any numerical
tolerance, so the bistability is genuine.

\paragraph{Spatial structure and mechanism of the hot-runaway state.}
The hot-runaway state [Fig.~\ref{fig:fig_manifold}(b)] is spatially
uniform with $J\!\simeq\!5.2$, $\theta\!\simeq\!12$,
$\varphi\!\simeq\!0.026$, and a thin reaction layer at the surface
where reactant is consumed almost entirely
($u_\mathrm{surf}\!\sim\!10^{-7}$).  The mechanism is diffusion-
limited combustion: the high-$\theta$ Arrhenius factor exceeds
$10^5$ and burns through any reactant in the boundary layer, but the
surface BC continues to deliver reactant at the diffusion-limited
rate $\Bi_c$.  The gel remains super-swollen because the high-$J$
Flory--Huggins minimum coexists with the collapsed minimum whenever
$\chi_\mathrm{eff}$ is sufficiently large; at $\varphi\!\sim\!0.026$
the $\chi_\mathrm{eff}\,\varphi^2$ penalty is negligible and the
elastic term $\Omega_e\,(J-1/J)$ becomes the dominant contribution to
the chemical potential.  The choice between the two
Flory--Huggins minima is set by the IC: trajectories landing in the
basin of the swollen branch reach hot-runaway, while trajectories
landing in the basin of the collapsed branch are ejected to the
LCST-front cycle by the surface BC and the diffusion-barrier
feedback.

\paragraph{Basin asymmetry and why the default sweep misses the
runaway.}
A refined IC sweep on a $15\times 15$ grid at the C-cell
[Fig.~\ref{fig:fig_manifold}(c)] places the basin boundary at
roughly $J_0\!\approx\!1.0$, $\theta_0\!\approx\!1$: only ICs in the
upper-left quadrant of $(J_0, \theta_0)$ space (cold, swollen) reach
the cycle.  ICs with any of (i)~$J_0\!\lesssim\!0.4$ (pre-collapsed),
(ii)~$\theta_0\!\gtrsim\!3$ (ignited), or (iii)~combinations
spanning the basin boundary fall into the runaway.  The basin of
the cycle covers
$\sim\!13\%$ of the IC grid (16 of the 119 successfully integrated
cells of an $11{\times}11$ grid spanning $J_0\in[0.20,1.30]$,
$\theta_0\in[0,10]$), a small fraction of the total. This asymmetry explains why the default
cold-IC sweep of Fig.~\ref{fig:fig4_phase} sees only the cycle and
never the runaway: the default cold IC sits well inside the cycle's
basin at every grid point of the parameter sweep.

\paragraph{IC-induced hysteresis: a falsifiable prediction.}
A 1D slice through the basin space gives a quantitative,
experimentally accessible signature of the bistability
[Fig.~\ref{fig:fig_manifold}(d)].  Holding all reaction and bath
parameters fixed at the C-cell value with the initial reactant
loading $u_0=0.5$, we vary the initial thermal pulse $\theta_0$ from
$0$ to $6$ for two seed conditions:
\begin{itemize}
\item ``cold'' seed ($J_0=1.30$): initially near the cold-swollen
  homogeneous state.  For $\theta_0$ below the threshold
  $\theta_0^{\!*}\!\approx\!2.4$ the trajectory remains in the
  cycle's basin and
  $\langle J\rangle_\mathrm{late}\!\approx\!2.0$;
  for $\theta_0$ above the threshold the seed crosses the basin
  boundary and
  $\langle J\rangle_\mathrm{late}\!\approx\!5.9$,
  indistinguishable from the runaway state.  The transition is
  sharp—the discontinuity in $\langle J\rangle_\mathrm{late}$ falls
  within a single $\Delta\theta_0=0.25$ sampling step.
\item ``pre-collapsed'' seed ($J_0=0.20$): initialized inside the
  runaway basin; for every $\theta_0$ in the scanned range the
  trajectory terminates in the runaway-type state with
  $\langle J\rangle_\mathrm{late}\!\approx\!4.2\text{--}5.5$,
  independent of $\theta_0$.
\end{itemize}
The split between the two terminal-state curves over the bistable
$\theta_0$ interval $[0, 2.4]$ is the hysteresis signature.
Experimentally, this would manifest as a sharp jump in long-time
gel volume (or temperature) when the initial perturbation magnitude
is varied across $\theta_0^{\!*}$, with no return jump on the
reverse sweep so long as the C-cell parameters are held fixed.
This is a quantitative prediction of the model that the linear
theory does not provide and that distinguishes spatially-extended
thermo-LCST gels from the well-mixed CSTR
analog~\cite{Uppal1974}.

\paragraph{Summary.}
The PDE attractor diagram is therefore richer than the single-IC
sweep of Fig.~\ref{fig:fig4_phase} suggests, and richer than 0D LSA
can predict.  The hot-runaway state is invisible to both linear
analysis and to default-IC numerical scans, and its appearance is
contingent on the structure of the basin of attraction in IC space.
Combined with the slow-manifold relaxation cycle of
Sec.~\ref{sec:lcst_front_mechanism}, this completes a four-attractor
inventory of the spatially extended problem (cold-swollen, LCST-
front cycle, frozen front, hot-runaway), whose existence and
stability are governed not by any single homogeneous-state
linearization but by the spatial coupling between the Flory--Huggins
free energy and the $\varphi$-dependent diffusion barrier.
